\section{The larger opportunity}\label{sec:horizon}
The motivating consequence is much larger than a better tariff. Many useful activities need several missing conditions to become an industry: a way to produce the result, a way to recognize it, a workable division of effort, a financing path and enough voluntary demand. Intelligence applied to that whole construction could make previously impractical cooperation repeatable.

The first consequential prototype should test a missing condition that matters beyond one convenient example. A small, decisive experiment can establish a large enabling relationship. Its ambition comes from the class of useful work it could open, not the size of a speculative market total or the number of agents assigned to it.

\subsection{Scientific capability as a shared productive service}
Suppose several research users need measurements that are individually too expensive to obtain. A construction might combine a reusable instrument, calibrated evidence and agreements that allocate access and support. The technical task is to obtain the required information reliably. The economic task is to make contributions, capacity and risk worthwhile for each participant. The strategic task includes alternative measurements, shared purchases, data derivation and competing suppliers.

The decisive first test would deliver a specific useful measurement under an independent acceptance condition, compare the strongest complete alternative, and then use the retained instrument on a distinct investigation. If the second investigation can already be done more cheaply with available equipment, the claimed enabling residual disappears. If it benefits only because the same data were handed to it, that is a content-reuse result. A stronger successor claim needs a new useful investigation and an attributable construction benefit.

\subsection{Industrial cooperation made possible by evidence}
Independent operators may possess complementary capacity while lacking a dependable way to coordinate its use. A productive arrangement could construct the required measurement, control interface and acceptance method, then allocate work and exposure under voluntary terms. A qualified intermediate output can permit earned progress payments; a meaningful measurement can make a contingent commitment credible.

The first test should measure the actual industrial outcome: conforming throughput, useful energy, reliable availability or another explicit beneficiary target. Settlement, maintenance, interruptions, shared resources and adverse conditions belong in the same arrangement. A simulated schedule does not operate equipment. A profitable model does not establish physical capacity. Moving to real devices requires the relevant engineering evidence and operating authority \cite{frontier}.

\subsection{A foundry for complete useful arrangements}
A broader foundry could construct and adapt the interfaces, verification routines and operating methods that connect specialist capabilities. The commercial product would be a useful accepted outcome or an independently usable productive capability. A simple existing tool may supply most of the answer; the valuable invention may be the small connector that removes a consequential obstacle.

Such a foundry should also search for ways to eliminate a dependency. If an inherited workflow requires reconstructing an unavailable state, the useful intervention may be a decision rule that needs less information. If repeated verification is expensive, a qualified production interface may make the relevant property directly checkable. Constructing every missing component of the original workflow is not always the best route to the authorized purpose \cite{ti,world,original}.

Each successful composite can become a candidate operator for further construction, with explicit applicability and invalidation conditions. A constructor of useful composites can itself become a retained object. This is the recursive ambition: the enterprise becomes better equipped to invent the conditions of subsequent progress. That last improvement must still earn the comparisons in Section~\ref{sec:inventor}.

\subsection{Create enough value to sustain the creator}
The enterprise needs a defensible share of the value it helps produce. Possible sources include superior accepted outcomes, reliable integration, calibrated evidence, shorter qualification, useful service continuity, scarce complementary assets and continuing productive invention. The relevant source depends on the market. A freely available implementation can make adoption and independent use easier, while a service or complementary productive asset remains costly to reproduce. None of those possibilities establishes willingness to pay on its own.

The created capability can also lower a rival's cost or let a buyer self-produce. This can be socially valuable while weakening the creator's margin. An invention program must therefore evaluate participant and competitive response alongside production. It should pursue a new instrument, a better delivery process, a new customer opportunity or retaining liquidity according to their evidence and complete consequences. Sustaining the enterprise is a real design problem; manufacturing dependence is not a substitute for creating value.

The laboratory's US\$3 billion objective concerns profit over a complete trailing 365-day window after modeled costs. It remains unachieved in the evidence reviewed here. Capital, financing, demand, price, repeat business, capacity, risk and competitive response all affect whether such a scale can be reached. There is no universal game-theoretic optimum without specifying those conditions, the feasible strategies and the objective. A finite optimum in the companion certifies only its supplied domain.

\section{A concrete engineering objective}\label{sec:objective}
An achieved autonomous economic-game inventor should be describable in operational terms. Given an authorized purpose and a declared environment, it constructs an executable arrangement, obtains the evidence needed to evaluate it, exposes consequential strategic and productive defects, and repairs or rejects the candidate. When a candidate is selected, it funds the required actions, delivers accepted work, settles the modeled or actual payments and carries the complete obligations. It retains useful means and tests their contribution to fresh successor work.

This objective has a meaningful bounded version. A system can complete that loop within an explicit finite language and environment. Its achievement should be reported with that boundary, including which primitive capabilities, type distributions, valuations, cost models and construction procedures were supplied. The broader version additionally requires discovery of materially new means or arrangements, robust advantage against capable alternatives, and evidence that retained procedures improve future invention after their acquisition cost.

For a consequential next experiment, select a beneficiary outcome with a clear missing condition and compare materially different routes: direct production or purchase, a new productive construction, and a different arrangement of evidence, commitments or work. Give the strongest complete alternative the same relevant freedom. Freeze the accepted outcome and finite comparison. Construct the smallest useful instrument that could decide the issue, then execute an entire funded episode if the evidence supports it. A failed test can identify the next construction or justify non-entry; it should not disappear from the cost history.

The publication and its companion support that work by clarifying one connection that a correct optimizer can miss: a productive invention changes the economic actions available to participants, and those actions change the resources available to the successor. The current finite engineering record already contains functioning mechanism construction, checked service production, simulated settlement and earned-funded continuation. It also contains conventional ties, adverse losses and a specific omitted-action failure. Those facts make the next construction more concrete.

\thesis{A lasting flywheel creates a possibility that works, leaves useful means behind, and earns the ability to construct what comes next.}

The research program is therefore both ambitious and falsifiable. Its unit of progress is a consequential accepted arrangement and a demonstrated successor relation. Its long horizon is an expanding repertoire of useful scientific, industrial and organizational capabilities. The work ahead is to make the machine increasingly able to discover and construct those arrangements, while participants remain free to choose and the evidence remains strong enough to tell us when it has succeeded.
